% 天体物理学（综述）
% license CCBYSA3
% type Wiki

本文根据 CC-BY-SA 协议转载翻译自维基百科\href{https://en.wikipedia.org/wiki/Astrophysics?utm_source=chatgpt.com}{相关文章}。


天体物理学是一门利用物理学与化学的方法和原理来研究天体及其现象的科学。\(^\text{[2]}\)作为该学科的奠基人之一，詹姆斯·基勒（James Keeler）曾指出，天体物理学“旨在探明天体的本质，而非它们在空间中的位置或运动——研究它们是什么，而非它们在哪里”，\(^\text{[3]}\)这一区别正体现于天体力学的研究范畴中。

天体物理学的研究对象包括太阳（太阳物理学）、其他恒星、星系、系外行星、星际介质以及宇宙微波背景辐射。\(^\text{[4][5]}\)科研人员在整个电磁波谱范围内观测这些天体的辐射，分析的物理特性包括光度、密度、温度及化学成分等。由于天体物理学的研究范围极其广泛，天体物理学家在研究中会综合运用多个物理学分支的概念与方法，包括经典力学、电磁学、统计力学、热力学、量子力学、相对论、核物理与粒子物理，以及原子与分子物理等。

在实践中，现代天文学研究通常涉及大量理论与观测物理学工作。天体物理学的主要研究领域包括：暗物质、暗能量、黑洞及其他天体的性质，以及宇宙的起源与最终命运。\(^\text{[4]}\)
理论天体物理学还研究多个主题，例如：太阳系的形成与演化、恒星动力学与演化、星系的形成与演化、磁流体力学、宇宙大尺度物质结构的形成、宇宙射线的起源、广义相对论与狭义相对论、量子宇宙学与物理宇宙学（即对宇宙最大尺度结构的物理研究），其中也包括弦论宇宙学与天体粒子物理等内容。
\subsection{历史}
天文学是一门古老的科学，长期以来与地球物理学相分离。在亚里士多德的宇宙观中，天体被认为是永恒不变的球体，其唯一的运动是匀速圆周运动；而地上世界则是生长与腐朽的领域，其中的自然运动沿直线进行，并在物体到达其目标时终止。因此，人们认为天界由一种与地球物质根本不同的物质构成——柏拉图认为那是火（Fire），而亚里士多德则称之为以太（Aether）。\(^\text{[6][7]}\)到了十七世纪，伽利略（Galileo）、笛卡尔（Descartes）和牛顿（Newton）等自然哲学家开始主张，天体与地上物体由相似的物质组成，并遵循相同的自然规律。\(^\text{[8][9][10][11]}\)然而，他们所面临的困难在于：当时尚未发明出能够证明这些观点的科学工具。\(^\text{[12]}\)

在十九世纪的大部分时间里，天文学研究的重点仍是例行性的观测工作，如测量天体的位置与计算其运动。\(^\text{[13][14]}\)一种新的天文学——很快被称为天体物理学（Astrophysics）——开始兴起。威廉·海德·沃拉斯顿（William Hyde Wollaston）与约瑟夫·冯·夫琅和费（Joseph von Fraunhofer）分别独立发现，当太阳光被分解时，其光谱中出现了大量暗线（即光强减弱或缺失的区域）。\(^\text{[15]}\)到 1860 年，物理学家古斯塔夫·基尔霍夫（Gustav Kirchhoff）与化学家罗伯特·本生（Robert Bunsen）证明，这些太阳光谱中的暗线与已知气体光谱中的亮线相对应，而每条特定的谱线都对应一种独特的化学元素。\(^\text{[16]}\)基尔霍夫推断，这些暗线是由太阳大气中化学元素吸收光线所造成的。\(^\text{[17]}\)由此，人们首次证明：太阳与恒星中存在的化学元素，也同样存在于地球上。

在进一步研究太阳与恒星光谱的人物中，诺曼·洛基尔（Norman Lockyer）尤为突出。1868 年，他在太阳光谱中发现了辐射线与暗线。与化学家爱德华·弗兰克兰（Edward Frankland）合作研究不同温度与压力下元素的光谱时，他注意到太阳光谱中一条黄色谱线无法对应任何已知元素。于是他提出，这条谱线代表一种新元素，并将其命名为“氦”（Helium），来源于希腊语“Helios”，意为“太阳之神”。\(^\text{[18][19]}\)

1885 年，爱德华·C·皮克林（Edward C. Pickering）在哈佛学院天文台发起了一项雄心勃勃的恒星光谱分类计划。他组织了一支女性计算员团队，其中包括威廉米娜·弗莱明（Williamina Fleming）、安东尼娅·莫里（Antonia Maury）与安妮·跳跃·坎农（Annie Jump Cannon）。她们根据照相底片上记录的光谱对恒星进行了系统分类。到 1890 年，团队已完成超过一万颗恒星的目录，并将它们划分为十三种光谱类型。根据皮克林的设想，到 1924 年，坎农将该目录扩展为九卷本、超过二十五万颗恒星，并发展出“哈佛分类系统”（Harvard Classification Scheme），该系统于 1922 年被国际天文学界正式采纳为全球标准。\(^\text{[20]}\)

1895 年，乔治·埃勒里·黑尔（George Ellery Hale）与詹姆斯·E·基勒（James E. Keeler）联合来自欧洲与美国的十位副主编，共同创办了《天体物理学期刊》（The Astrophysical Journal: An International Review of Spectroscopy and Astronomical Physics）。\(^\text{[21][22]}\)该期刊旨在填补天文学与物理学期刊之间的学术空白，为发表以下研究提供平台：关于光谱仪在天文学中的应用；与天体物理密切相关的实验室研究，包括金属与气体光谱的波长测定、辐射与吸收实验；关于太阳、月球、行星、彗星、流星与星云的理论研究；以及望远镜与实验室仪器的改进与设计。\(^\text{[21]}\)

约在 1920 年，随着赫罗图（Hertzsprung–Russell Diagram）的发现——这一图表至今仍是恒星分类与演化研究的基础——阿瑟·爱丁顿（Arthur Eddington）在论文《恒星的内部结构》（The Internal Constitution of the Stars）中预见了恒星内部核聚变过程的存在与机制。\(^\text{[23][24]}\)当时，恒星能量的来源仍是一个完全的谜团。爱丁顿正确地推测，这一能量来源于氢核聚变为氦的过程，并根据爱因斯坦方程 $E = mc^2$ 解释了能量的释放。这一推断极具前瞻性，因为当时科学界尚未发现核聚变、热核能，甚至尚未明确恒星主要由氢构成（参见“金属丰度”metallicity）。\(^\text{[25]}\)

1925 年，塞西莉亚·海伦娜·佩恩（Cecilia Helena Payne，后为 Cecilia Payne-Gaposchkin）在拉德克利夫学院（Radcliffe College）完成了一篇具有深远影响的博士论文。她将萨哈电离理论（Saha’s Ionization Theory）应用于恒星大气，首次将恒星光谱类型与其温度建立了系统联系。\(^\text{[26]}\)更重要的是，她发现恒星的主要成分是氢与氦，而非与地球相似的化学组成。尽管爱丁顿对此表示支持，但这一结论过于出人意料，以至于她的论文审阅者（包括亨利·诺里斯·罗素 Henry Norris Russell）劝说她在发表前修改了结论。然而，后续研究最终证实了她的发现。\(^\text{[27][28]}\)

到 20 世纪末，对天体光谱的研究已扩展至覆盖从无线电波到可见光、X 射线与伽马射线的广泛波段。\(^\text{[29]}\)进入 21 世纪，观测范围进一步拓展，纳入了基于引力波探测的新型天文观测手段。
\subsection{观测天体物理学}
\begin{figure}[ht]
\centering
\includegraphics[width=6cm]{./figures/fa6be55c6141f1c4.png}
\caption{} \label{fig_AsYs_1}
\end{figure}
观测天文学是天文学的一个分支，主要致力于观测数据的记录与解释；与之相对的是理论天体物理学（Theoretical Astrophysics），后者主要研究物理模型可被测量的推论。  
观测天文学的核心实践是使用望远镜及其他天文仪器来观测天体。

大多数天体物理观测都是通过电磁波谱进行的。
\begin{itemize}
\item 射电天文学研究波长大于数毫米的辐射。典型的研究对象包括：由星际气体与尘埃云等低温天体所发射的射电波；宇宙微波背景辐射——即大爆炸（Big Bang）红移后的光；以及脉冲星（Pulsar），其最初便是在微波频段被探测到的。研究此类射电波需要使用口径巨大的射电望远镜。
\item 红外天文学研究波长过长而无法被肉眼看到、但又短于射电波的辐射。  
红外观测通常使用与可见光望远镜类似的仪器进行。比恒星温度更低的天体（例如行星）通常通过红外波段进行研究。
\item 光学天文学是最早发展的天文学类型。常用仪器包括装配有电荷耦合器件（CCD）或光谱仪（Spectroscope）的望远镜。地球大气会在一定程度上干扰光学观测，因此人们使用自适应光学系统（Adaptive Optics）与空间望远镜（Space Telescopes）来获得尽可能高质量的图像。在这一波段范围内，恒星极为明亮，同时可通过化学光谱分析恒星、星系与星云的化学组成。
\item 紫外线、X 射线与伽马射线天文学研究极高能的天体物理过程，例如双星脉冲星系统（Binary Pulsars）、黑洞（Black Holes）、磁星（Magnetars）等。这类辐射无法有效穿透地球大气，因此科学家采用两种方法观测该波段的电磁辐射：空间望远镜与地面成像大气契伦科夫望远镜（Imaging Air Cherenkov Telescopes, IACT）。前者的代表性天文台包括 RXTE、钱德拉 X 射线天文台（Chandra X-ray Observatory）与康普顿伽马射线天文台（Compton Gamma Ray Observatory）；后者的代表性设施包括高能立体系统（High Energy Stereoscopic System, H.E.S.S.）与 MAGIC 望远镜。
\end{itemize}
除电磁辐射之外，从地球上能够观测到的、起源于遥远宇宙的信号极为有限。目前已建成少数引力波观测台，但引力波（Gravitational Waves）的探测极其困难。此外，人们还建造了中微子观测台（Neutrino Observatories），主要用于研究太阳。由极高能粒子组成的宇宙射线（Cosmic Rays）也可以通过其撞击地球大气层时产生的效应进行观测。

观测的时间尺度也可能存在显著差异。大多数光学观测的时间范围为数分钟至数小时，因此变化速度更快的天体现象往往难以直接观测。然而，一些天体具有跨越数百年甚至数千年的历史观测记录。另一方面，射电观测既可以分析毫秒量级的事件（如毫秒脉冲星，Millisecond Pulsars），也可以结合多年数据进行分析（例如脉冲星减速研究）。这些不同时间尺度的观测所得信息差异极大。

太阳的研究在观测天体物理学中占据特殊地位。由于其他恒星距离极为遥远，太阳是唯一能够以无与伦比的细节进行直接观测的恒星。理解太阳的结构与活动机制，有助于我们理解其他恒星的物理特性与演化规律。

关于恒星如何变化，即恒星演化（Stellar Evolution）的研究，通常通过将不同类型的恒星放置在赫罗图（Hertzsprung–Russell Diagram）上加以建模。该图可视为描述恒星从诞生到毁灭整个生命周期状态的示意图。
\subsection{理论天体物理学}
理论天体物理学家使用多种研究工具，包括解析模型（例如利用多重指数模型Polytropes来近似恒星的行为）与计算数值模拟。  
两者各有优点：解析模型通常更能揭示物理过程的本质；而数值模型则能发现那些否则无法直接观测到的现象与效应。\(^\text{[30][31]}\)

天体物理学的理论研究者致力于建立理论模型，并推导出这些模型的观测结果，以便观测者能够寻找数据来验证或反驳模型，或在多个相互竞争的模型之间作出选择。理论学家同时也会根据新的观测数据来修正模型。若出现不一致之处，通常会尽量以最小改动来调整模型以符合数据。然而，若长期积累的大量数据均与模型相矛盾，往往会导致该模型被彻底放弃。

理论天体物理学的研究主题包括：恒星动力学与演化；星系的形成与演化；磁流体力学（Magnetohydrodynamics）；宇宙大尺度物质结构的形成；宇宙射线的起源；广义相对论与物理宇宙学（Physical Cosmology），其中还包括弦论宇宙学（String Cosmology）与天体粒子物理（Astroparticle Physics）。相对论性天体物理学（Relativistic Astrophysics）作为分析引力在大尺度结构中重要作用的理论工具，构成了黑洞物理学与引力波研究的基础。

目前在天体物理学中被广泛接受并深入研究的理论与模型，包括大爆炸理论（Big Bang）、宇宙暴胀（Cosmic Inflation）、暗物质（Dark Matter）、暗能量（Dark Energy）以及各种基本物理理论，它们现已共同被纳入$\Lambda$CDM 模型之中。
\subsection{大众化与传播}
天体物理学的起源可追溯至十七世纪统一物理学的兴起，当时学者首次提出相同的自然法则同时适用于天界与尘世。\(^\text{[11]}\)早期既精通物理又研究天文学的科学家，为现代天体物理学奠定了坚实基础。

在现代，学生们之所以持续对天体物理学产生兴趣，部分原因在于其被英国皇家天文学会（Royal Astronomical Society）及多位杰出科学传播者所推广——其中包括著名教授劳伦斯·克劳斯（Lawrence Krauss）、苏布拉马尼扬·钱德拉塞卡（Subrahmanyan Chandrasekhar）、史蒂芬·霍金（Stephen Hawking）、于贝尔·里夫斯（Hubert Reeves）、卡尔·萨根（Carl Sagan）以及帕特里克·穆尔（Patrick Moore）。这些学者与机构在不同时期的努力，持续吸引着年轻一代投入天体物理学的学习与研究。\(^\text{[32][33][34]}\)

美国情景喜剧《生活大爆炸》（The Big Bang Theory）也在大众文化中普及了天体物理学这一领域，剧中更邀请了多位著名科学家客串演出，例如史蒂芬·霍金（Stephen Hawking）与尼尔·德格拉斯·泰森（Neil deGrasse Tyson）。
\subsection{参见}
\begin{itemize}
  \item \textbf{天体化学（Astrochemistry）} —— 研究宇宙中分子的存在与化学反应的学科。
  \item \textbf{天文台（Astronomical Observatories）} —— 用于观测与研究天体的设施。
  \item \textbf{天文光谱学（Astronomical Spectroscopy）} —— 测量天体电磁辐射的科学方法。
  \item \textbf{天体粒子物理学（Astroparticle Physics）} —— 粒子物理学的一个分支，研究宇宙中的高能粒子。
  \item \textbf{引力波天文学（Gravitational-wave Astronomy）} —— 利用引力波观测宇宙的天文学分支。
  \item \textbf{赫罗图（Hertzsprung–Russell Diagram）} —— 展示恒星光度与光谱类型关系的散点图。
  \item \textbf{高能天文学（High-energy Astronomy）} —— 研究发射高能电磁辐射的天体。
  \item \textbf{天体物理学重要出版物（Important Publications in Astrophysics）}
  \item \textbf{天文学家列表（List of Astronomers）} —— 包括天体物理学家在内的学者名录。
  \item \textbf{中微子天文学（Neutrino Astronomy）} —— 通过观测低质量恒星粒子研究宇宙的领域。
  \item \textbf{引力物理与相对论发展时间线（Timeline of Gravitational Physics and Relativity）}
  \item \textbf{关于星系、星系团与宇宙大尺度结构的知识发展时间线（Timeline of Knowledge about Galaxies, Clusters of Galaxies, and Large-scale Structure）}
  \item \textbf{白矮星、中子星与超新星研究时间线（Timeline of White Dwarfs, Neutron Stars, and Supernovae）} —— 记录这些领域的知识发展与观测历史的时间顺序列表。
\end{itemize}
\subsection{参考文献}
\begin{enumerate}
  \item Maoz, Dan（2016）。Astrophysics in a Nutshell（天体物理学概要）。普林斯顿大学出版社，第272页。ISBN~978-1400881178。
  \item “astrophysics（天体物理学）”。Merriam–Webster 出版公司。原始版本存档于 2011 年 6 月 10 日。检索日期：2011 年 5 月 22 日。
  \item Keeler, James~E.（1897 年 11 月）。“天体物理学研究的重要性及其与其他物理科学的关系”（The Importance of Astrophysical Research and the Relation of Astrophysics to the Other Physical Sciences）。  
  \textit{The Astrophysical Journal}，6（4）：271–288。Bibcode:~1897ApJ.....6..271K。doi:~10.1086/140401。PMID~17796068。
  \item “Focus Areas – NASA Science（NASA 科学研究重点领域）”。nasa.gov。原始存档于 2017 年 5 月 16 日。检索日期：2017 年 7 月 12 日。
  \item “astronomy（天文学）”。Encyclopædia Britannica（大英百科全书）。2023 年 5 月 29 日。
  \item Lloyd, G.~E.~R.（1968）。Aristotle: The Growth and Structure of His Thought（亚里士多德思想的发展与结构）。剑桥：剑桥大学出版社，第134–135页。ISBN~978-0-521-09456-6。
  \item Cornford, Francis~MacDonald（约 1957）[1937]。Plato's Cosmology: The Timaeus of Plato translated, with a running commentary（柏拉图的宇宙论：〈蒂迈欧篇〉英译与评注）。印第安纳波利斯：Bobbs Merrill Co.，第118页。
  \item Galilei, Galileo（1989），Van~Helden, Albert（编）。Sidereus Nuncius or The Sidereal Messenger（星际信使）。芝加哥：芝加哥大学出版社，第21、47页。ISBN~978-0-226-27903-9。
  \item Slowik, Edward（2013）[2005]。“Descartes’ Physics（笛卡尔的物理学）”。载于 Stanford Encyclopedia of Philosophy（斯坦福哲学百科全书）。检索日期：2015 年 7 月 18 日。
  \item Westfall, Richard~S.（1983）。Never at Rest: A Biography of Isaac Newton（永不止息：艾萨克·牛顿传）。剑桥：剑桥大学出版社（1980 出版），第731–732页。ISBN~978-0-521-27435-7。
  \item Burtt, Edwin~Arthur（2003）[首版 1924]。The Metaphysical Foundations of Modern Science（现代科学的形而上学基础）。第二修订版。纽约 Mineola：Dover Publications，第30、41、241–242页。ISBN~978-0-486-42551-1。
  \item Kvasz, Ladislav（2013）。“Galileo, Descartes, and Newton – Founders of the Language of Physics（伽利略、笛卡尔与牛顿——物理学语言的奠基者）”。Acta Physica Slovaca（斯洛伐克物理学会学报）。捷克科学院哲学研究所出版。检索日期：2015 年 7 月 18 日。
  \item Case, Stephen（2015）。“‘Land-marks of the Universe’: John Herschel against the Background of Positional Astronomy（〈宇宙的地标〉：约翰·赫歇尔与位置天文学的背景）”。Annals of Science（科学年鉴），72（4）：417–434。Bibcode:~2015AnSci..72..417C。doi:~10.1080/00033790.2015.1034588。PMID~26221834。S2CID~205397708。  
  文中指出：19 世纪早期的大多数天文学家对恒星的物理性质并无兴趣，恒星被视为用于建立精确参考背景以描绘太阳、月球与行星运动的“标尺”，主要服务于地面应用。
  \item Donnelly, Kevin（2014 年 9 月）。《科学的乏味：十九世纪的位置天文学》（On the Boredom of Science: Positional Astronomy in the Nineteenth Century）。The British Journal for the History of Science（英国科学史杂志），47（3）：479–503。doi:~10.1017/S0007087413000915。S2CID~146382057。
  \item Hearnshaw, J.~B.（1986）。The Analysis of Starlight（星光分析）。剑桥：剑桥大学出版社，第23–29页。ISBN~978-0-521-39916-6。
  \item Kirchhoff, Gustav（1860）。《论夫琅禾费线》（Ueber die Fraunhofer'schen Linien）。  
  \textitAnnalen der Physik（物理学年刊），185（1）：148–150。Bibcode:~1860AnP...185..148K。doi:~10.1002/andp.18601850115。
  \item Kirchhoff, Gustav（1860）。《论物体热与光的发射能力与吸收能力之间的关系》（Ueber das Verhältniss zwischen dem Emissionsvermögen und dem Absorptionsvermögen der Körper für Wärme und Licht）。Annalen der Physik（物理学年刊），185（2）：275–301。Bibcode:~1860AnP...185..275K。doi:~10.1002/andp.18601850205。
  \item Cortie, A.~L.（1921）。《诺曼·洛克耶爵士（1836–1920）》（Sir Norman Lockyer, 1836–1920）。  
  \textit{The Astrophysical Journal（天体物理学杂志）}，53：233–248。Bibcode:~1921ApJ....53..233C。doi:~10.1086/142602。
  \item Jensen, William~B.（2004）。《为什么氦元素以“-ium”结尾》（Why Helium Ends in ``-ium''）(PDF)。Journal of Chemical Education（化学教育期刊），81（7）：944–945。Bibcode:~2004JChEd..81..944J。doi:~10.1021/ed081p944。
  \item Hetherington, Norriss~S.; McCray, W.~Patrick; Weart, Spencer~R.（编）。Spectroscopy and the Birth of Astrophysics（光谱学与天体物理学的诞生）。美国物理学会（American Institute of Physics），物理史中心（Center for the History of Physics）。原始文献存档于 2015 年 9 月 7 日。检索日期：2015 年 7 月 19 日。
  \item Hale, George~Ellery（1895）。《天体物理学杂志》（The Astrophysical Journal）。The Astrophysical Journal（天体物理学杂志），1（1）：80–84。Bibcode:~1895ApJ.....1...80H。doi:~10.1086/140011。
  \item The Astrophysical Journal（天体物理学杂志），1（1）。
  \item Eddington, A.~S.（1920 年 10 月）。《恒星的内部结构》（The Internal Constitution of the Stars）。The Scientific Monthly（科学月刊）}，11（4）：297–303。Bibcode:~1920Sci....52..233E。doi:~10.1126/science.52.1341.233。JSTOR~6491。PMID~17747682。
  \item Eddington, A.~S.（1916）。《论恒星的辐射平衡》（On the Radiative Equilibrium of the Stars）。Monthly Notices of the Royal Astronomical Society（皇家天文学会月刊）}，77：16–35。Bibcode:~1916MNRAS..77...16E。doi:~10.1093/mnras/77.1.16。
  \item McCracken, Garry；Stott, Peter（2013）。McCracken, Garry；Stott, Peter（编）。Fusion（聚变）（第二版）。波士顿：Academic Press，第13页doi:~10.1016/b978-0-12-384656-3.00002-7。ISBN~978-0-12-384656-3。文中指出：Eddington 认识到，当四个氢原子结合成一个氦原子时会出现质量损失。爱因斯坦的质能等价原理直接引出了这样一个假设：这正是恒星能量产生的长期寻找的机制！这是一个极具洞察力的推测，尤其值得注意的是，当时原子核结构和反应机制尚未被完全理解。
  \item Payne, C.~H.（1925）。Stellar Atmospheres: A Contribution to the Observational Study of High Temperature in the Reversing Layers of Stars（恒星大气层：对恒星反转层高温的观测性研究的贡献）（博士论文）。美国马萨诸塞州剑桥市：Radcliffe College。Bibcode:~1925PhDT.........6P。
  \item Haramundanis, Katherine（2007）。《佩恩–加波什金［佩恩］，塞西莉亚·海伦娜》（Payne–Gaposchkin [Payne], Cecilia Helena）。载于：Hockey, Thomas；Trimble, Virginia；Williams, Thomas~R.（编），Biographical Encyclopedia of Astronomers（天文学家传记百科全书）}。纽约：Springer出版社，第876–878页。ISBN~978-0-387-30400-7。检索日期：2015 年 7 月 19 日。
  \item Soter, Steven；Tyson, Neil~deGrasse（2000）。《塞西莉亚·佩恩与恒星的组成》（Cecilia Payne and the Composition of the Stars）American Museum of Natural History（美国自然历史博物馆）。
  \item Biermann, Peter~L.; Falcke, Heino（1998）。《天体物理学前沿：研讨会总结》（Frontiers of Astrophysics: Workshop Summary）。载于：Panvini, Robert~S.; Weiler, Thomas~J.（编），Fundamental Particles and Interactions: Frontiers in Contemporary Physics, an International Lecture and Workshop Series（基本粒子与相互作用：当代物理学前沿——国际讲座与研讨系列）。AIP Conference Proceedings（美国物理学会会议录），第423卷。美国物理学会，第236–248页。arXiv:~astro-ph/9711066。Bibcode:~1998AIPC..423..236B。doi:~10.1063/1.55085。ISBN~1-56396-725-1。

  \item Roth, H.（1932）。《\textit{缓慢收缩或膨胀的流体球及其稳定性}》（A Slowly Contracting or Expanding Fluid Sphere and Its Stability）。  
  \textit{Physical Review（物理评论）}，39（3）：525–529。Bibcode:~1932PhRv...39..525R。doi:~10.1103/PhysRev.39.525。

  \item Eddington, A.~S.（1988）[1926]。《\textit{恒星的内部结构}》（Internal Constitution of the Stars）。  
  \textit{Science（科学）}，52（1341）。纽约：剑桥大学出版社，第233–240页。  
  Bibcode:~1920Sci....52..233E。doi:~10.1126/science.52.1341.233。ISBN~978-0-521-33708-3。PMID~17747682。

  \item Manley, D.~Mark（2012）。《\textit{著名的天文学家与天体物理学家}》（Famous Astronomers and Astrophysicists）。  
  美国肯特州立大学（Kent State University）。检索日期：2015 年 7 月 17 日。

  \item The~science.ca~team（2015）。《\textit{于贝尔·里夫斯——天文学、天体物理学与空间科学}》（Hubert Reeves – Astronomy, Astrophysics and Space Science）。  
  GCS~Research~Society。检索日期：2015 年 7 月 17 日。

  \item “\textit{Neil deGrasse Tyson（尼尔·德格拉斯·泰森）}”。海登天文馆（Hayden Planetarium）。2015 年。检索日期：2015 年 7 月 17 日。
\end{enumerate}

\end{enumerate}



